\documentclass[11pt]{article}
\usepackage[margin=2.5cm]{geometry}
\usepackage{amsmath,amssymb,amsthm}
\usepackage[hidelinks]{hyperref}
\newcommand{\Q}{\mathbb{Q}}
\newcommand{\Z}{\mathbb{Z}}
\newcommand{\N}{\mathbb{N}}
\newcommand{\R}{\mathbb{R}}
\newcommand{\p}{\mathfrak{p}}
\newcommand{\m}{\mathfrak{m}}
\newcommand{\cR}{\mathcal{R}}
\newcommand{\cF}{\mathcal{F}}
\newcommand{\inv}{\operatorname{inv}}
\newcommand{\maxinv}{\operatorname{maxinv}}
\newcommand{\gr}{\operatorname{gr}}
\numberwithin{equation}{section}
\theoremstyle{plain}
\newtheorem{theorem}{Theorem}[section]
\newtheorem{proposition}[theorem]{Proposition}
\newtheorem{lemma}[theorem]{Lemma}
\newtheorem{corollary}[theorem]{Corollary}
\theoremstyle{definition}
\newtheorem{definition}[theorem]{Definition}
\newtheorem{example}[theorem]{Example}
\theoremstyle{remark}
\newtheorem*{comment}{Comment}
\newtheorem*{remark}{Remark}
\title{A constructive proof that a Bézout domain\\ need not be an elementary divisor domain}
\author{Claude (Anthropic)}
\date{7 October 2026}
\begin{document}
\maketitle

\begin{abstract}
We give a constructive proof of the theorem of Hägg and Mörtberg, with GPT-6 Astra: there is a
Bézout domain containing $\Q[x,y]$ over which the matrix
$\bigl(\begin{smallmatrix}1+x&y\\y&1-x\end{smallmatrix}\bigr)$ has no Smith normal form. The proof
is written in the style of~\cite{LQ}. The only principles used beyond intuitionistic logic are
unique choice and Noetherianity in the sense of Richman and Seidenberg. Krull
dimension, maximal ideals and compactness are not used.

This note is a constructive version, written by Claude, of the proof in~\cite{HM}
(\url{https://arxiv.org/abs/2609.35229}). It has not been reviewed by Hägg and Mörtberg. The
accompanying Lean formalization, a fork of theirs, is at
\url{https://github.com/coquand/bezout-constructive}.
\end{abstract}

\tableofcontents

\section{Statement and plan}

\begin{theorem}[Hägg--Mörtberg with GPT-6 Astra~\cite{HM}]\label{t:main}
There is a Bézout domain $R\supseteq\Q[x,y]$ over which the matrix
\[
M=\begin{pmatrix}1+x&y\\y&1-x\end{pmatrix}
\]
is not equivalent to a diagonal matrix. In particular $M$ has no Smith normal form.
\end{theorem}

The ring is constructed as an increasing union $R=\bigcup_n A_n$ of finitely presented
$\Q$-algebras with $A_0=\Q[x,y]$. Every pair of elements of $R$ is enumerated. At stage $n$ one
pair $(a,b)$ of $A_n$ is written $(a,b)=c\,(a',b')$ with $\gcd(a',b')=1$, and $A_{n+1}$ is a ring in
which $(a',b')$ becomes principal. This step is a \emph{principalisation}, done by repeating two
moves (divisorial and torsor) until an invariant can no longer decrease.

The obstruction is real and topological. Put $\Delta=x^2+y^2-1$. Then
\[
\det M=-\Delta,\qquad M^2-2M=\Delta I,
\]
and on the circle $\Delta=0$, written $(x,y)=(a^2-b^2,2ab)$ with $a^2+b^2=1$, one has
$M=2\,{}^t(a,b)\,(a,b)$. So $\frac12 M$ projects onto the line of the half-angle, a Möbius
band. The final argument (Section~\ref{s:end}) shows that a diagonal form over some $A_N$ would give
a nowhere-vanishing section of this band along a closed path of real points.

\paragraph{What has to be made constructive.} The plan above is that of~\cite{HM}. Every step
uses classical notions: maximal ideals, the maximum of an invariant over them, minimal primes,
Krull dimension, compactness of real points, and factoriality. We replace them as follows.
\begin{itemize}
\item Factoriality becomes \emph{discrete gcd domain} (Section~\ref{s:gcd}).
\item Smoothness and dimension are \emph{certificates}: explicit Jacobian covers
  (Section~\ref{s:cert}).
\item Maximal ideals become \emph{points} with values in number fields (Section~\ref{s:points}).
\item Noetherianity in the sense of Richman--Seidenberg (Section~\ref{s:noeth}) gives
  termination of searches without a priori bounds.
\item Compactness becomes a \emph{chain lifting property} for finite chains of real points
  (Section~\ref{s:chains}).
\end{itemize}

\section{Conventions}\label{s:conv}

We work in Bishop-style constructive mathematics, as in~\cite{LQ}.
\begin{itemize}
\item A ring is \emph{discrete} if equality is decidable. An ideal is \emph{detachable} if
  membership is decidable. A ring is \emph{strongly discrete} if every finitely generated ideal is
  detachable. Finitely presented $\Q$-algebras are strongly discrete and coherent (Gröbner bases,
  \cite[Chapter~VII]{LQ}).
\item A \emph{discrete field} is a discrete ring in which every element is $0$ or invertible.
  A \emph{number field} is a finite-dimensional $\Q$-algebra that is a discrete field.
\item Real numbers are Cauchy sequences with modulus. We use: if $a<b$ then $x<b$ or $a<x$
  (cotransitivity); inverses of positive reals; $n$-th roots of positive reals; and the fixed point of
  a contraction with rate $\frac12$. We never test a real number for zero.
\end{itemize}

\paragraph{Unique choice.} If for every $x$ there is a unique $y$ with $P(x,y)$, then there is a
function $f$ with $P(x,f(x))$. In Bishop's mathematics this is part of the notion of function. The
construction of the tower is a recursion in which each stage is the unique object satisfying a
decidable specification (canonical choices, Section~\ref{s:locus}). So it is a function by unique
choice, and no countable or dependent choice is used.

\section{Noetherianity and Krull's intersection theorem}\label{s:noeth}

We use Noetherianity in the sense of Richman and Seidenberg~\cite[II-3.2]{LQ}: every ascending
sequence of finitely generated ideals has two equal consecutive terms. Equivalently, for every
sequence $(a_k)$ there is $j$ with $a_j\in\langle a_0,\dots,a_{j-1}\rangle$.

\begin{proposition}\label{p:hilbert}
Finitely generated $\Q$-algebras, their quotients and localisations, and the Rees algebras of
their finitely generated ideals (Definition~\ref{d:rees}) are Noetherian.
\end{proposition}

\begin{proof}
Hilbert's basis theorem in the constructive form of~\cite{CP}. To pass to quotients and
localisations without countable choice, one uses an inductive formulation of Noetherianity; see
Appendix~\ref{a:bar}.
\end{proof}

The proofs below all follow one pattern. A process produces a sequence step by step. Noetherianity
gives an index where the new term lies in the ideal of the earlier ones, and this contradicts the
hypothesis that keeps the process running. So the process stops, although no bound is known in
advance.

\begin{lemma}[order along a non-unit]\label{l:sorder}
Let $D$ be a Noetherian domain with decidable divisibility by $s$, where $s\in D$ is neither $0$ nor
a unit. Every $a\ne0$
can be written $a=s^ka'$ with $s\nmid a'$, and $k$ is computed.
\end{lemma}

\begin{proof}
Put $a_0=a$. Put $a_{i+1}=a_i/s$ if $s\mid a_i$ (decidable), and $a_{i+1}=a_i$ otherwise. By
Noetherianity there is $j\ge1$ with $a_j\in\langle a_0,\dots,a_{j-1}\rangle$. If $s$ divided
$a_0,\dots,a_{j-1}$, this ideal would be $\langle a_{j-1}\rangle$ and $a_{j-1}=sa_j$. Then
$a_j=c\,s\,a_j$, so $(1-cs)a_j=0$. Since $s$ is not a unit, $1-cs\ne0$, so $a_j=0$ and hence
$a=0$, a contradiction. So the division stops at some $i<j$, which is found by a finite check.
\end{proof}

\begin{definition}[{\cite[XII-2, before Fact~2.5]{LQ}}]\label{d:rees}
Let $\mathfrak{a}$ be an ideal of a ring $B$ and $t$ an indeterminate. The \emph{Rees algebra} of
$\mathfrak{a}$ is the subalgebra
\[
B[\mathfrak{a}t]=B\oplus\mathfrak{a}t\oplus\mathfrak{a}^2t^2\oplus\cdots\ \subseteq B[t].
\]
So a polynomial $\sum g_jt^j$ lies in $B[\mathfrak{a}t]$ iff $g_j\in\mathfrak{a}^j$ for every $j$.
If $\mathfrak{a}=\langle a_1,\dots,a_k\rangle$, then $B[\mathfrak{a}t]$ is generated by $a_1t,\dots,a_kt$,
so it is a quotient of $B[T_1,\dots,T_k]$.
\end{definition}

\begin{lemma}[Krull step, after Perdry~\cite{Pm}]\label{l:krull}
Let $\mathfrak{a}$ be an ideal of a ring $B$. Suppose that in the Rees ring $B[\mathfrak{a}t]$ a
relation
\[
b\,t^{m+1}=\sum_{i\le m}g_i\,b\,t^i,\qquad g_i\in B[\mathfrak{a}t],
\]
holds. Then $b\in b\,\mathfrak{a}$. If moreover $1+\mathfrak{a}\subseteq B^\times$, then $b=0$.
\end{lemma}

\begin{proof}
Compare the coefficients of $t^{m+1}$. The coefficient of $t^{m+1-i}$ in $g_i$ lies in
$\mathfrak{a}^{m+1-i}\subseteq\mathfrak{a}$, so $b=cb$ with $c\in\mathfrak{a}$, and then $(1-c)b=0$.
\end{proof}

\begin{corollary}[Krull's intersection theorem]\label{c:krull}
Let $B$ be a local ring with maximal ideal $\m$, whose Rees algebra $B[\m t]$ is Noetherian, and
in which membership in each $\m^k$ is decidable. For $b\in B$ one computes either $k$ with
$b\notin\m^k$, or a proof that $b=0$.
\end{corollary}

\begin{proof}
Test $b\in\m^k$ for $k=0,1,\dots$. As long as the tests succeed, $bt^k\in B[\m t]$. Noetherianity
applied to the sequence $(bt^k)$ gives the relation of Lemma~\ref{l:krull}, and then $b=0$, because
$1+\m\subseteq B^\times$.
\end{proof}

\section{Discrete gcd domains}\label{s:gcd}

\begin{definition}
A \emph{discrete gcd domain} is a discrete domain with decidable divisibility in which every two
elements have a gcd.
\end{definition}

This replaces factoriality. It is all we need to write any pair as $(a,b)=c(a',b')$ with $a',b'$
without common non-unit factor.

\begin{lemma}\label{l:poly}
$\Q[x,y]$ is a discrete gcd domain. If $D$ is a discrete gcd domain, so are $D[X]$ and every
localisation of $D$ at a multiplicative set with decidable membership.
\end{lemma}

\begin{proof}
Gauss's lemma, and gcds over the discrete field of fractions~\cite[III]{LQ}. For $\Q[x,y]$ one can
also use Kronecker factorisation.
\end{proof}

\begin{lemma}[Nagata at a prime element]\label{l:nagata}
Let $D$ be a strongly discrete Noetherian domain and $s\in D$ a prime element, that is,
$D/sD$ is a domain. If $D[1/s]$ is a discrete gcd domain, so is $D$.
\end{lemma}

\begin{proof}
Let $a,b\ne0$. By Lemma~\ref{l:sorder} write $a=s^ka'$ and $b=s^lb'$ with $s\nmid a',b'$. Let
$g\in D$ represent a gcd of $a',b'$ in $D[1/s]$; dividing out powers of $s$ (Lemma~\ref{l:sorder}
again), we may assume $s\nmid g$.
\begin{itemize}
\item $g\mid a'$ in $D$. Indeed $s^Na'=gx$ for some $N$ and $x\in D$. Since $s$ is prime and
  $s\nmid g$, $s^N\mid x$, so $a'=g\,(x/s^N)$.
\item A common divisor $c\in D$ of $a',b'$ divides $g$ in $D$. It divides $g$ in $D[1/s]$, and
  $s\nmid c$ (because $s\nmid a'$), so the same argument applies.
\end{itemize}
Hence $s^{\min(k,l)}g$ is a gcd of $a,b$ in $D$. The factor $s^{\min(k,l)}$ is correct because $s$
is prime and $s\nmid a'$, $s\nmid b'$.
\end{proof}

\begin{lemma}[generic linear form]\label{l:generic}
Let $B$ be a discrete gcd domain and $c_0,\dots,c_\ell\in B$ with $\gcd(c)=1$. Over
$B'=B[\lambda_0,\dots,\lambda_\ell]$ the element $\pi'=\sum\lambda_ic_i$ is prime.
\end{lemma}

\begin{proof}
$B'$ is a discrete gcd domain. $\pi'$ has degree $1$ in the $\lambda$'s and content $\gcd(c)=1$. In a
factorisation one factor has degree $0$ and divides every coefficient, so it is a unit, and $\pi'$
is irreducible. In a discrete gcd domain an irreducible element $\pi$ is prime. If $\pi\mid ab$ and
$\pi\nmid a$, then $\gcd(\pi,a)=1$, hence $\pi\mid b$.
\end{proof}

For a ring $B$ and $c=(c_0,\dots,c_\ell)$ put
\[
J_B(c)=B[\sigma_0,\dots,\sigma_\ell]\Big/\Big(\sum c_i\sigma_i-1\Big),
\]
the \emph{forcing algebra} of $c$.

\begin{lemma}\label{l:jdomain}
If $D$ is a domain and some $c_i\ne0$, then $J_D(c)$ is a domain.
\end{lemma}

\begin{proof}
The polynomial $f=\sum c_i\sigma_i-1$ has content $(1)$. By the Dedekind--Mertens lemma, $J_D(c)$ is
torsion-free over $D$, so it embeds in $K[\sigma]/(f)$ with $K=\operatorname{Frac}D$. There $c_i$
is invertible, $\sigma_i$ can be solved, and what remains is a polynomial ring over $K$.
\end{proof}

\begin{proposition}\label{p:jgcd}
Let $B$ be a finitely presented $\Q$-algebra that is a discrete gcd domain, and let
$c_0,\dots,c_\ell\in B$ have no common non-unit divisor, with $c_0\ne0$. Then $J_B(c)$ is a discrete
gcd domain.
\end{proposition}

\begin{proof}
Let $B'=B[\lambda_0,\dots,\lambda_\ell]$ and $\pi'=\sum\lambda_ic_i$, which is prime by
Lemma~\ref{l:generic}. The substitution $\tau_i=\sigma_i+\lambda_i\sigma'$ gives
\[
J_B(c)[\lambda,\sigma']\;\cong\;E=B'[\tau,\sigma']\Big/\Big(\sum c_i\tau_i-\pi'\sigma'-1\Big).
\]
\begin{itemize}
\item $\pi'$ is prime in $E$: $E/\pi'E\cong J_{B'/\pi'}(\bar c)[\sigma']$, which is a domain by
  Lemma~\ref{l:jdomain}, because $B'/\pi'$ is a domain and $\bar c_0\ne0$ ($\pi'\nmid c_0$ for
  degree reasons).
\item Inverting $\pi'$ solves $\sigma'$, so $E[1/\pi']$ is a polynomial ring over $B'[1/\pi']$, hence
  a discrete gcd domain (Lemma~\ref{l:poly}).
\item $E$ is finitely presented, hence Noetherian with decidable divisibility (Gröbner bases).
\end{itemize}
So $E$ is a discrete gcd domain by Lemma~\ref{l:nagata}. Gcds of elements of $J_B(c)$ computed in the
polynomial ring $E\cong J_B(c)[\lambda,\sigma']$ have degree $0$, so they are gcds in $J_B(c)$.
\end{proof}
In the torsor step (Section~\ref{s:steps}), Lemma~\ref{l:nagata} is applied once more, at the
element $s$.

\section{Smoothness certificates}\label{s:cert}

Let $A=\Q[Y_1,\dots,Y_N]/(G)$ be finitely presented.

\begin{definition}\label{d:cert}
A \emph{certificate of dimension $n$} for $A$ consists of the following data.
\begin{itemize}
\item Polynomials $h_1,\dots,h_r$ and $u_1,\dots,u_r$, and an integer $C>0$, with
  $C-\sum u_kh_k\in(G)$. So the $h_k$ are comaximal in $A$.
\item For each $k$, polynomials $F=(F_1,\dots,F_N)$ with $F_1,\dots,F_{N-n}\in(G)$, together with
  $U$ and exponents $e,f$, such that modulo $(G)$
  \[
  \det\Bigl(\frac{\partial F}{\partial Y}\Bigr)U\equiv h_k^{\,e}
  \qquad\text{and}\qquad
  h_k^{\,f}\,G\subseteq(F_1,\dots,F_{N-n}).
  \]
\end{itemize}
The last $n$ components $F_{N-n+1},\dots,F_N$ are the \emph{coordinates} of the piece.
\end{definition}

So on $D(h_k)$ the ideal $(G)$ is generated by $N-n$ equations, and their Jacobian, completed by
the coordinates, is invertible.

\begin{lemma}[standard form of a piece]\label{l:etale}
Write $h=h_k$. Then
\[
A[1/h]\;\cong\;\Q[Y,Z_1,\dots,Z_n,t]\big/\bigl(F_1,\dots,F_{N-n},\ F_{N-n+j}-Z_j,\ ht-1\bigr),
\]
and the Jacobian of this presentation with respect to $(Y,t)$ is a unit. Hence $A[1/h]$ is étale
over $\Q[Z_1,\dots,Z_n]$, and $\Omega_{A[1/h]/\Q}$ is free with basis $dZ_1,\dots,dZ_n$.
\end{lemma}

\begin{proof}
The isomorphism is given by two explicit inverse maps. No kernel has to be computed. The Jacobian
is $\det(\partial F/\partial Y)\cdot h$, which is a unit modulo the relations.
\end{proof}

From Lemma~\ref{l:etale} we get explicit local coordinates, explicit dual derivations
$\partial/\partial Z_j$, and Newton lifting of solutions at points. These are all the smoothness
properties we use; no general smoothness criterion is needed.

\begin{remark}[smooth versus smooth of constant dimension]
A certificate says more than ``$A$ is smooth over $\Q$'': it fixes the rank $n$ of
$\Omega_{A/\Q}$ on every piece. For one equation $F$, the condition that $\{F=0\}$ is a smooth
hypersurface,
\[
1\in\Bigl(F,\frac{\partial F}{\partial Y_1},\dots,\frac{\partial F}{\partial Y_N}\Bigr),
\]
is a certificate of dimension $N-1$, with $h_k=\partial F/\partial Y_k$. Smoothness alone allows
the rank to vary (for instance $X^2=I$ in $M_n$, which has isolated points). In our tower all rings
are discrete domains, whose only idempotents are $0$ and $1$, so the rank is automatically
constant. The certificate records it.
\end{remark}

\begin{proposition}[propagation]\label{p:prop}
\begin{enumerate}
\item If $B$ has a certificate of dimension $n$ and $c_0,\dots,c_r\in B$, then $J_B(c)$ has a certificate of dimension $n+r$.
\item If $A$ has a certificate of dimension $n$, the extended Rees algebra of a weighted centre
  (defined in Section~\ref{s:steps}) has a certificate of dimension $n+1$.
\end{enumerate}
\end{proposition}

\begin{proof}
1. The $c_i$ are comaximal in $J=J_B(c)$, since $\sum c_i\sigma_i=1$ there. On $D(c_i)$ one solves
$\sigma_i$: $J[1/c_i]\cong B[1/c_i][\sigma_j:j\ne i]$. The pieces are the products $c_ih_k$, and the
coordinates are those of $B$ together with the $\sigma_j$, $j\ne i$. In the language of the remark,
$\sum c_i\sigma_i-1$ is a smooth hypersurface relative to $B$, because the derivatives
$\partial/\partial\sigma_i=c_i$ generate $(1)$ in $J$.
2. On a chart where the centre is $(x_1,\dots,x_k)$ with integer weights $w_i$,
\[
\cR\;\cong\;A[s,u_1,\dots,u_k]\big/(x_i-s^{w_i}u_i),
\]
and the equations $x_i-s^{w_i}u_i$ are solved for $x_i$, which are coordinates of $A$. The new
coordinates are those of $A$ with $x_i$ replaced by $u_i$, plus $s$.
\end{proof}

The certificate gives a dimension without Krull dimension. All rings of the tower have
dimension $n\ge2$. The dimension $n$ bounds the length of the invariants: $k\le n$
(Section~\ref{s:inv}).

\section{Points}\label{s:points}

\begin{definition}
A \emph{point} of a finitely presented $\Q$-algebra $A$ is a $\Q$-algebra map $q:A\to L$, where
$L$ is a number field (Section~\ref{s:conv}). Its kernel $\ker q$ is a detachable prime, and
its residue field $\kappa(q)\subseteq L$ is a discrete field.
\end{definition}

A point of $\Q[Y]/(G)$ is a solution of $G=0$ in $L^N$. In the classical proof, every statement
``for every maximal ideal $\m\supseteq I$'' becomes ``for every point $q$ with $I\subseteq\ker q$''.
Classically the two are equivalent by the Nullstellensatz. Constructively we use the following.

\begin{proposition}[Nullstellensatz]\label{p:null}
Let $\p$ be a finitely generated prime of $A$ and $a\notin\p$. There is a point $q$ with
$\p\subseteq\ker q$ and $q(a)\ne0$. More generally, if $I$ is a finitely generated ideal with
$1\notin I$, there is a point with $I\subseteq\ker q$.
\end{proposition}

\begin{proof}
$A/I$ is a nonzero finitely presented $\Q$-algebra (decided by Gröbner bases). Noether
normalisation and the primitive element theorem produce a quotient that is a finite field
extension~\cite[VII]{LQ}. For the first claim, apply this to $(A/\p)[1/a]$.
\end{proof}

\begin{lemma}[transfer]\label{l:transfer}
\begin{enumerate}
\item Membership in $\ker q$ is decidable. Hence a negative statement of the form $x\notin\ker q$
  may be proved by contradiction.
\item If $\varphi:A\to U$ is a $\Q$-algebra map and $Q$ a point of $U$, then $Q\circ\varphi$ is a
  point of $A$. This replaces the classical fact that the contraction of a maximal ideal of a
  finitely generated algebra is maximal.
\item (Local--global) If $g_1,\dots,g_m$ are comaximal and a property of finitely presented modules
  or ideals holds on each $D(g_i)$, it holds on $A$~\cite[II]{LQ}.
\end{enumerate}
\end{lemma}

\begin{proposition}[covering the zero set]\label{p:cover}
Let $I$ be a finitely generated ideal of $A$. Suppose that for every point $q$ with
$I\subseteq\ker q$ we can compute $g_q\notin\ker q$. Then finitely many points $q_1,\dots,q_m$
satisfy $1\in I+(g_{q_1},\dots,g_{q_m})$.
\end{proposition}

\begin{proof}
Build the list $g_{q_1},g_{q_2},\dots$. If $1\in I+(g_{q_1},\dots,g_{q_j})$ (decidable) stop.
Otherwise Proposition~\ref{p:null} gives a point $q_{j+1}$ of $V(I+(g_{q_1},\dots,g_{q_j}))$. Then
$g_{q_{j+1}}\notin I+(g_{q_1},\dots,g_{q_j})$, because $q_{j+1}$ kills the latter. Continue the sequence by $1$
after a stop, and apply the Noetherianity of $A/I$ (Proposition~\ref{p:hilbert}). It gives $j$ with
$g_{q_{j+1}}\in I+(g_{q_1},\dots,g_{q_j})$. By the last observation, the process has stopped by
step~$j$.
\end{proof}

This is our substitute for quasi-compactness of $\operatorname{Spec}A$.

\section{The invariant at a point}\label{s:inv}

\subsection{Centred charts and Taylor expansion}

Let $q$ be a point of $A$, lying in a piece $D(h)$ of the certificate, and let $Z_1,\dots,Z_n$ be the
coordinates of the piece. After a translation by the values $q(Z_j)$ (and, when $\kappa(q)\ne\Q$,
after adjoining the residue field), the coordinates vanish at $q$; such a chart is \emph{centred}
at $q$. Let $A_q$ be the localisation at $\ker q$, and $\m_q$ its maximal ideal. Lemma~\ref{l:etale}
gives the \emph{Taylor map}
\[
\tau:A_q\to\kappa(q)[[Z_1,\dots,Z_n]],
\]
obtained by Newton lifting. Its coefficients are computed. By Corollary~\ref{c:krull} it is
injective. Moreover $\m_q^k$ consists exactly of the elements whose expansion has no terms of degree
$<k$.

\begin{lemma}[order at a point]\label{l:ord}
For $f\in A_q$ we can decide whether $f=0$, and if $f\ne0$ compute $\mathrm{ord}_q(f)$, the largest
$k$ with $f\in\m_q^k$. The same holds for a finitely generated ideal $I$, as the minimum over its
generators.
\end{lemma}

\begin{proof}
Test the homogeneous parts of $\tau(f)$ of degree $0,1,2,\dots$; each test is a finite computation
over $\kappa(q)$. Corollary~\ref{c:krull} says that this search ends, either at a nonzero part or
with a proof that $f=0$. No bound is known in advance, and no Markov principle is used.
\end{proof}

Weighted orders work in the same way. For rational weights $e_1,\dots,e_k\ge0$ on some coordinates,
$f$ lies in the filtration level $\cF_a$ if and only if $\tau(f)$ has no monomial $Z^\beta$ with
$\sum e_i\beta_i<a$. In particular, the associated graded ring $\bigoplus\cF_a/\cF_{>a}$ embeds in
weighted-homogeneous series over $\kappa(q)$, so it is a domain.

\subsection{The invariant}

At a point $q\supseteq I$ we use the invariant of Abramovich--Temkin--Włodarczyk~\cite{ATW}, in the
differential presentation of Brais~\cite{Brais} (``Method~1''). It is a finite sequence
\[
\inv_q(I)=(a_1,\dots,a_k),\qquad k\le n,
\]
of positive rationals, computed as follows.
\begin{itemize}
\item $a_1=\mathrm{ord}_q(I)$ (Lemma~\ref{l:ord}).
\item A \emph{maximal contact} element is a derivative of order $a_1-1$ of a generator of order
  $a_1$. It has order $1$ at $q$, and it is found by a finite search.
\item The next entries come from the coefficient ideal on the maximal contact hypersurface. Each
  is a minimum over a finite, explicitly bounded box of multi-indices. It has the form $b/D!$ with
  $b\in\N_{>0}$, where $D$ is the product of the denominators of the previous entries.
\end{itemize}
Equivalently one may record the weights $w_i=1/a_i$. Then $\maxinv$ corresponds to the
lexicographically \emph{least} weight vector.
Every step is a finite computation over the discrete field $\kappa(q)$. The same data give a
\emph{marked centre} at $q$: coordinates $x_1,\dots,x_k$ and weights $w_1,\dots,w_k$. Its
admissibility ($I\subseteq\cF_{d}$ for the appropriate $d$) is decided coefficientwise.

\begin{lemma}[well-ordering]\label{l:wo}
Order invariants lexicographically, with the convention of~\cite{ATW} for sequences of different
lengths. Restricted to sequences of length $\le n$ that arise as invariants, this order is well
founded: well-founded induction and recursion are valid on it
(Appendix~\ref{a:bar}).
\end{lemma}

\begin{proof}
Induction on the length. Once $a_1,\dots,a_{i-1}$ are fixed, the next entry is $b/D!$ with $D$
fixed by the previous entries. So it ranges over $\frac1{D!}\N$, which is well founded by induction
on $\N$. The lexicographic product of well-founded orders
is well founded, by nested induction.
\end{proof}

\begin{proposition}[upper semicontinuity]\label{p:usc}
For every point $q\supseteq I$ there is $g\notin\ker q$ such that $\inv_{q'}(I)\preceq\inv_q(I)$
for every point $q'\supseteq I$ of $D(g)$; the computation runs at every point of $D(g)$.
Moreover, on $D(g)$ the points with
$\inv_{q'}(I)=\inv_q(I)$ are those of $V(x_1,\dots,x_k)$, for the marked centre at $q$.
\end{proposition}

\begin{proof}[Sketch]
All the data of the computation at $q$ (a nonzero coefficient, a unit derivative, an admissible
marked centre) are elements that do not vanish at $q$. Let $g$ be their product. At a point $q'$ of
$D(g)$ the same data are still available, so the computation at $q'$ cannot produce a larger
invariant. It produces the same invariant exactly when the marked centre vanishes at $q'$. Each
inequality is a statement about non-membership in $\ker q'$, which is decidable
(Lemma~\ref{l:transfer}).
\end{proof}

\section{The maximal locus and its components}\label{s:locus}

\begin{proposition}\label{p:maxinv}
For a finitely generated ideal $I$ with $0\ne I\ne(1)$, one can compute $\maxinv(I)$, a point $q_0$ where it is attained,
and the property
\[
\inv_q(I)\preceq\maxinv(I)\quad\text{for every point }q\supseteq I .
\]
\end{proposition}

\begin{proof}
Combine Proposition~\ref{p:usc} with Proposition~\ref{p:cover}. This gives finitely many points
$q_1,\dots,q_m$ whose neighbourhoods $D(g_{q_j})$ cover $V(I)$. Take the maximum of the
$\inv_{q_j}(I)$. Every point $q\supseteq I$ lies in some $D(g_{q_j})$, decidably, so
$\inv_q(I)\preceq\inv_{q_j}(I)$.
\end{proof}

The \emph{locus ideal} $L$ is obtained by gluing: on the $D(g_{q_j})$ with $\inv_{q_j}=\maxinv$ it is
the marked centre ideal, and elsewhere it is $(1)$. Its points are exactly the points where the
maximum is attained. Since the centre is given by part of a coordinate system, $V(L)$ is smooth.

\paragraph{Components.} The \emph{components} are the minimal primes of $\sqrt L$. They are
computed with Perdry's strongly Noetherian theory~\cite{P}: a primality test by induction on the
number of variables, and a tree of prime-or-split decisions whose finiteness again comes from Noetherianity. Different components $\p\ne\p'$ are comaximal. Indeed, $1\in\p+\p'$ is decidable. If it fails,
Proposition~\ref{p:null} gives a point $q$ of $V(\p+\p')$. Near $q$ the locus is $V(x_1,\dots,x_k)$
for a centred chart, and these coordinates generate a prime. Both $\p$ and $\p'$ are minimal over the
locus, so both are equal to this prime, and $\p=\p'$.

\paragraph{Canonical choice.} Each finitely presented ring of the tower comes with an explicit
enumeration. Order the finite lists of generators by first occurrence in that enumeration. The
component treated next is the one whose canonical list of generators is least. Together with the
canonical point $q_0$ (least in the same sense), all choices are canonical, so the step is a
function by unique choice.

\begin{remark}
One could try a dynamical treatment: keep a radical centre and split it only when needed. This does
not work for the tower. With $r\ge2$ components, the forcing algebra has class group
$\Z^r/\Z(1,\dots,1)\ne0$. A split discovered later would therefore invalidate a stage that has
already been used.
\end{remark}

\section{The principalisation step}\label{s:steps}

Let $I=(a',b')$ in a stage $A$ with a certificate of dimension $n$, let $\maxinv(I)$ have first
entry $a_1$, and let $\p$ be the chosen component, with weighted centre $x_1^{1/w_1},\dots,
x_k^{1/w_k}$ along $\p$.

\subsection{Divisorial step ($k=1$)}

Then $\p$ is principal. Let $\pi$ be a gcd of a generating family of $\p$. At one point of $V(\p)$,
$\p$ is locally principal, and gcds of families are preserved by localisation, so $\p=(\pi)$ and
$\pi$ is prime.
The weight is $1/a$, and $I=\pi^aI_1$. The new ideal is $I_1$, in the same ring.

\begin{lemma}
No point of $V(\pi)$ contains $I_1$. So the component $\p$ disappears from the locus of $I_1$, and the
invariant does not increase elsewhere.
\end{lemma}

\begin{proof}
Let $q\supseteq I_1$ with $\pi\in\ker q$. Then $I=\pi^aI_1\subseteq\m_q^{a+1}$, so the first weight
at $q$ would be at most $1/(a+1)$, contradicting the maximality of $\maxinv(I)$, whose first entry is $a$. Membership in $\ker q$ is
decidable, so this negative argument is valid (Lemma~\ref{l:transfer}).
\end{proof}

\subsection{Torsor step ($k\ge2$)}

We need two variants of Definition~\ref{d:rees}: the filtration is weighted, and negative degrees
are allowed. Let $\cF_j$ ($j\in\Z$) be the \emph{weighted filtration} of the centre. On a chart
where the centre has coordinates $x_1,\dots,x_k$ and integer weights $w_1,\dots,w_k$, $\cF_j$ is the
ideal generated by the monomials $x^\beta$ with $\sum w_i\beta_i\ge j$; in particular $\cF_j=A$
for $j\le0$, and $\cF_i\cF_j\subseteq\cF_{i+j}$. These local ideals glue. The \emph{extended Rees
algebra} is the subalgebra
\[
\cR=\bigoplus_{j\in\Z}\cF_jT^j\ \subseteq A[T,T^{-1}],\qquad s=T^{-1}\in\cR .
\]
On the chart, $\cR=A[s,\,x_1T^{w_1},\dots,x_kT^{w_k}]$. Putting $u_i=x_iT^{w_i}$ gives
$x_i=s^{w_i}u_i$, which is the presentation used in Proposition~\ref{p:prop}. Moreover
$\cR/s\cR\cong\bigoplus_{j\ge0}\cF_j/\cF_{j+1}=\gr\cF$, the associated graded ring, and
$\cR[1/s]=A[T,T^{-1}]$. Choose $\pi\in\cF_1\setminus\cF_2$, and let $y=(y_0,\dots,y_\ell)=(\pi T,h_1,\dots,h_\ell)$, where the $h_i$
generate $\cR_+$ (in practice $x_iT^{w_i}$). Put
\[
U=J_{\cR}(y)=\cR[\sigma_0,\dots,\sigma_\ell]\Big/\Big(\sum y_i\sigma_i-1\Big).
\]
The weights of the centre are $v_i=w_i/d$ for a common denominator $d$, and admissibility gives
$I\subseteq\cF_d$. Every $f\in I$ satisfies $f=s^d(fT^d)$, so
\[
IU=s^d\,I_1U,\qquad I_1=(fT^d : f\in I)\subseteq\cR,
\]
where $I_1$ is the \emph{weak transform} of $I$.
By Proposition~\ref{p:prop}, $U$ has a certificate of dimension $n+1+\ell$.

\begin{proposition}\label{p:sprime}
$s$ is a prime element of $U$, and $U$ is a discrete gcd domain.
\end{proposition}

\begin{proof}
$U/sU$ is the forcing algebra of the images of the $y_i$ over $\cR/s\cR=\gr\cF$, and the image of
$\pi T$ is nonzero. By Lemma~\ref{l:jdomain} it is enough that $\gr\cF$ is a domain. That is checked at \emph{one} point $q_0$ of $V(\p)$: by Section~\ref{s:inv},
$\gr\cF$ embeds into weighted-homogeneous series over $\kappa(q_0)$. A statement $z\notin\cF_{j+1}$ is
transferred from $q_0$ to the whole component by decidability. Then $U[1/s]=J_{B}(y)$ with $B=A[T,T^{-1}]$. Here $y_0=\pi T\ne0$, and the $y_i$ have no
common non-unit divisor in $B$: up to units they are generators of $\cF_1\supseteq\p$, and $k\ge2$.
So $U[1/s]$ is a discrete gcd domain (Proposition~\ref{p:jgcd}). Lemma~\ref{l:nagata}
concludes.
\end{proof}

\begin{proposition}[the invariant drops]\label{p:drop}
For every point $Q$ of $U$ with $I_1\subseteq\ker Q$, $\inv_Q(I_1)\prec\maxinv(I)$, or $Q$ lies over
a component $\p'\ne\p$, where the invariant is unchanged.
\end{proposition}

\begin{proof}[Sketch]
Let $q=Q|_A$ (Lemma~\ref{l:transfer}). There are two cases.
\begin{itemize}
\item \emph{Away from the centre.} If $q\notin V(\p)$, then $U$ is locally a polynomial ring over
  $A$, and nothing changes. The other components lift to the primes $\p'U$, which are comaximal
  with $\p U$. So the locus ideal of $I_1$ is $\bigcap_{\p'\ne\p}\p'U$, which has one component
  less.
\item \emph{Over the centre}, with $s\in\ker Q$; then $\cR_+\not\subseteq\ker Q$, since the $y_i$
  generate $(1)$ in $U$. Let $Q_0$ be the \emph{vertex point} of $\cR$: the point
  $\sum f_jT^j\mapsto q_0(f_0)$, which lies in $V(\cR_+)$. At $Q_0$ the weak transform $I_1$ has
  invariant $\maxinv(I)$, attained along the Rees centre, which lies in $V(\cR_+)$. On a
  neighbourhood $D(G)$ of $Q_0$, Proposition~\ref{p:usc} bounds the invariant by $\maxinv(I)$, with
  equality only on the Rees centre. Hence $\inv_Q(I_1)\prec\maxinv(I)$ as soon as $Q$ restricts to a
  point of $D(G)$. To arrange this we use the $\Q^\times$-action $T\mapsto\mu T$, which preserves
  $I_1$ up to a unit. $Q(\sigma_\mu G)$ is a polynomial in $\mu$ with coefficients in the number
  field of $Q$, of some degree $e$, and its constant term $Q(G_0)$ is nonzero. So one of
  $\mu=1,\dots,e+1$ is not a root, and it is found by the zero test.\qedhere
\end{itemize}
\end{proof}

\begin{theorem}[principalisation]\label{t:princ}
Let $A$ be a stage with a certificate, $(a',b')$ without common non-unit factor. A finite sequence of
divisorial and torsor steps produces a stage $A'\supseteq A$ with a certificate, a discrete gcd
domain, and $g\in A'$ with $(a',b')A'=gA'$. The construction is a function of $(A,a',b')$.
\end{theorem}

\begin{proof}
The pair $(\maxinv(I),\ \text{number of components})$ decreases lexicographically at each step: the
divisorial step and the torsor step remove the chosen component, without increasing the invariant
elsewhere. By Lemma~\ref{l:wo} the order is well founded, so the steps are defined by well-founded
recursion. The recursion stops when $I=(1)$. The accumulated factor
$\pi^a$ or $s^d$ of each step generates the original ideal. Every choice is canonical
(Section~\ref{s:locus}), so unique choice turns the recursion into data.
\end{proof}

\subsection{The tower}

Let $A_0=\Q[x,y]$, with the trivial certificate of dimension $2$. Enumerate pairs so that every pair
of elements of every $A_n$ is treated at some later stage. At stage $n$, write the pair as
$c\,(a',b')$ using the gcd of $A_n$, and apply Theorem~\ref{t:princ}. Then $R=\bigcup_nA_n$ is a
domain. Every pair $(a,b)$ becomes principal at some stage, and stays principal, so $R$ is a
Bézout domain. The maps $A_n\to A_{n+1}$ are injective, since the forcing algebra is faithful when
$\mathrm{Ann}(c)=0$~\cite[II.2.14]{LQ} and the Rees algebra contains $A$.

\section{Chain lifting}\label{s:chains}

A \emph{real point} of $A_n=\Q[Y]/(G)$ is a solution of $G=0$ in $\R^N$; let $X_n\subseteq\R^N$ be the
set of real points. A \emph{closed $\varepsilon$-chain} is a finite sequence $P_0,\dots,P_m$ with
$|P_{i+1}-P_i|<\varepsilon$ and $P_m=P_0$.

\begin{definition}\label{d:chain}
A map $\varphi:X'\to X$ has the \emph{chain lifting property} if for every bound $C$ there is a
bound $C'$ such that for every $\varepsilon>0$ there is $\delta>0$ with the following property.
Every closed $\delta$-chain $P_0,\dots,P_m$ in $X\cap[-C,C]^N$ has a closed $\varepsilon$-chain
$P'_0,\dots,P'_{m'}$ in $X'\cap[-C',C']^{N'}$ and a map $\rho:\{0,\dots,m'\}\to\{0,\dots,m\}$ with
\[
\rho(0)=0,\quad\rho(m')=m,\quad\rho(j+1)-\rho(j)\in\{0,1\},\quad
|\varphi(P'_j)-P_{\rho(j)}|<\varepsilon .
\]
\end{definition}

This is a finitary form of ``$\varphi$ is a proper surjection with connected fibres''. It mentions
only finitely many points at a time.

\begin{lemma}\label{l:compose}
If $\varphi$ and $\psi$ have the chain lifting property and $\varphi$ is Lipschitz on bounded sets,
then $\varphi\circ\psi$ has the chain lifting property.
\end{lemma}

\begin{proof}
Lift first along $\varphi$ with a small $\varepsilon_1$, then along $\psi$, and compose the
reindexings. A Lipschitz constant $K$ of $\varphi$ on the relevant bounded set turns $\varepsilon$
errors upstairs into $K\varepsilon$ errors downstairs. The parameters are chosen in the order $C$,
$C'$, $C''$, and then $\varepsilon$, $\delta$.
\end{proof}

The projections of the tower are polynomial maps, so they are Lipschitz on bounded sets. It
remains to treat one step. The divisorial step does not change the ring, and a principalisation is
a finite composite of steps.

\begin{proposition}\label{p:onestep}
For a torsor step $A\to U$, the induced map $X_U\to X_A$ has the chain lifting property.
\end{proposition}

\begin{proof}[Sketch]
Let $\nu=\sum_i g_i^{E_i}$ be the weighted distance to the centre, with exponents making the terms
weighted-homogeneous of a common degree $D$. Then two-sided bounds $c\,r^D\le\nu\le c'r^D$ hold in
chart coordinates.
\begin{itemize}
\item \emph{Away from the centre}, where $\nu>\eta$, the lift is unique and explicit:
  $s=\nu^{1/2D}$ (a positive root), and $\sigma$ is determined by a fixed formula. Lifting point
  by point gives the chain.
\item \emph{Near the centre}, the fibre over a point is a weighted sphere $\{\varphi(u)=1\}$ of
  dimension $k-1\ge1$, which is connected. Consecutive lifted points are joined inside the fibre by
  relay paths made of three linear segments in chart coordinates. Charts are changed through
  $\nu$.
\item The analytic input is a polynomial inverse function theorem, whose radius is uniform in the
  base point on bounded sets. It is proved by a Kantorovich--Newton step, i.e.\ the fixed point of
  a contraction with rate $\frac12$. It uses only inverses and roots of positive reals; there is no
  real division by a number not known to be positive, and no Markov principle.\qedhere
\end{itemize}
\end{proof}

\begin{corollary}\label{c:tower}
For every $N$, the map $X_{A_N}\to X_{A_0}=\R^2$ has the chain lifting property.
\end{corollary}

\section{Proof of Theorem~\ref{t:main}}\label{s:end}

Suppose $PMQ=\operatorname{diag}(e_1,e_2)$ with $P,Q$ invertible over $R$. Over a Bézout domain,
$\operatorname{diag}(e_1,e_2)$ is equivalent to $\operatorname{diag}(g,e_1e_2/g)$ with
$g=\gcd(e_1,e_2)$, so we may assume $e_1\mid e_2$. All entries lie in some $A_N$. Then $e_1$ divides
the entries of $M$, which generate $(1)$ because $(1+x)+(1-x)=2$. So $e_1$ is a unit. With $(p,q)$ the first row of $P$, $(u,v)$ the first column of
$Q$ and $\mu=e_1^{-1}$,
\begin{equation}\label{e:id}
\mu\,(p,q)\,M\binom uv=1\quad\text{in }A_N .
\end{equation}

\paragraph{Half-angles.} For $-1\le t\le1$ put
\[
(a(t),b(t))=\Bigl(\frac{1-t^2}{1+t^2},\ \frac{2t}{1+t^2}\Bigr),\qquad
\gamma(t)=(a^2-b^2,\ 2ab)\in\{\Delta=0\}.
\]
As $t$ runs from $-1$ to $1$, $\gamma$ runs once around the circle, from $(-1,0)$ back to $(-1,0)$,
while $(a,b)$ goes from $(0,-1)$ to $(0,1)$. On the circle $M=2\,{}^t(a,b)(a,b)$, so
\[
M\binom uv=2\,l\binom ab,\qquad l=ua+vb .
\]

\paragraph{A lower bound.} On a bounded region the functions $p,q,u,v,\mu$ (polynomials in the
coordinates of $A_N$) are bounded, and $u,v$ are Lipschitz. By~\eqref{e:id}, at a real point of
$A_N$ lying over a point of the circle, $2\,\mu\,l\,(pa+qb)=1$, so $|l|\ge 2c$ for some $c>0$. At a
real point $P'$ lying $\varepsilon$-close to $\gamma(t)$, the same bound holds up to an error that
tends to $0$ with $\varepsilon$. So $|l(P',t)|\ge c$ once $\varepsilon$ is small enough.

\paragraph{The chain.} Choose $-1=t_0<\dots<t_m=1$ so fine that $P_i=\gamma(t_i)$ is a closed
$\delta$-chain. Lift it by Corollary~\ref{c:tower} to a closed $\varepsilon$-chain $P'_j$ with
reindexing $\rho$. Put $l_j=u(P'_j)a(t_{\rho(j)})+v(P'_j)b(t_{\rho(j)})$. For small $\varepsilon$
and $\delta$:
\begin{itemize}
\item $|l_j|\ge c$ for all $j$;
\item $|l_{j+1}-l_j|<c$, because $u,v$ are Lipschitz, $t\mapsto(a,b)(t)$ is $4$-Lipschitz, and
  $\rho(j+1)-\rho(j)\in\{0,1\}$.
\end{itemize}
By cotransitivity, applied to $-c/2<c/2$, each $l_j$ satisfies $l_j>c/2$ or $l_j<-c/2$, hence
$l_j\ge c$ or $l_j\le-c$. A sign change between consecutive values would give
$|l_{j+1}-l_j|\ge2c$, so two consecutive values have the same sign. So $l_{m'}$ and $l_0$ have the same sign. But $P'_{m'}=P'_0$, and
$(a,b)(t_m)=-(a,b)(t_0)$, so $l_{m'}=-l_0$. This is a contradiction. Hence $M$ has no diagonal form
over $R$.\qed

\begin{comment}
This last step is the constructive content of the intermediate value theorem used classically.
It is where the ordered field $\R$ enters. Over a field in which $-1$ is a square, the circle ring
is a principal ideal domain and there is no obstruction.
\end{comment}

\section{What is used}

\begin{itemize}
\item Intuitionistic logic, unique choice, and inductive definitions (Appendix~\ref{a:bar}).
\item Constructive algebra from~\cite{LQ}: Gröbner bases, the Nullstellensatz, local--global
  principles, gcds over a discrete field, Kronecker factorisation, primitive elements.
\item The Hilbert basis theorem in inductive form~\cite{CP}, and Perdry's theory of minimal
  primes~\cite{P}.
\item On $\R$: cotransitivity, inverses and roots of positive numbers, and the fixed point of a
  contraction with rate $\frac12$.
\end{itemize}
No maximal ideals, Krull dimension, compactness, Markov principle or countable choice are used.


\appendix
\section{Inductive definitions and bar induction}\label{a:bar}

This appendix explains the inductive notions behind Proposition~\ref{p:hilbert} and
Lemma~\ref{l:wo}. They are not needed to follow the main text.

\paragraph{Inductive definitions.} An inductive definition introduces a property by rules: the
property is the smallest one closed under the rules. The natural numbers are generated by $0$ and
$n\mapsto n+1$, and induction on $\N$ says exactly that every property closed under these rules
holds for all $n$. The same is true for any inductive definition. This is \emph{induction on the
definition}. In type theory inductive definitions are primitive. In Bishop's style they are
accepted as generalised inductive definitions, as~\cite{LQ} does for the Krull dimension
($\mathrm{Hdim}$).

\begin{definition}[accessibility]
Let $\prec$ be a relation on a set $E$. The property $\mathrm{Acc}(x)$ is defined inductively by one
rule: if $\mathrm{Acc}(y)$ for all $y\prec x$, then $\mathrm{Acc}(x)$. The relation is \emph{well
founded} if $\mathrm{Acc}(x)$ holds for every $x$.
\end{definition}

Induction on $\mathrm{Acc}$ is well-founded induction, and it justifies well-founded recursion: a
function may be defined at $x$ from its values at all $y\prec x$. This is how the principalisation
is defined (Theorem~\ref{t:princ}). The lexicographic product of two well-founded relations is
well founded, by nested induction; this gives Lemma~\ref{l:wo}.

\paragraph{Bars.} Let $A$ be a commutative ring. Lists are written with the newest entry first.

\begin{definition}[inductive Noetherianity]\label{d:bar}
The property $\mathrm{Bar}(l)$ of a list $l$ of elements of $A$ is defined inductively by two rules.
\begin{enumerate}
\item If $a\in\langle l\rangle$, then $\mathrm{Bar}(a::l)$.
\item If $\mathrm{Bar}(a::l)$ for every $a\in A$, then $\mathrm{Bar}(l)$.
\end{enumerate}
$A$ is \emph{inductively Noetherian} if $\mathrm{Bar}$ holds for the empty list.
\end{definition}

Read $\mathrm{Bar}(l)$ as ``however $l$ is continued, it eventually gets an entry in the ideal of
the older ones''. This form is due to Richman, Seidenberg, Jacobsson--Löfwall and
Coquand--Persson~\cite{CP}.

\begin{lemma}\label{l:seq}
An inductively Noetherian ring is Noetherian in the sense of Richman--Seidenberg.
\end{lemma}

\begin{proof}
By induction on $\mathrm{Bar}(l)$, prove: ``every sequence beginning with $l$ has a term in the
ideal of the earlier terms''. Under rule~1, the newest entry of $l$ is such a term. Under rule~2,
apply the induction hypothesis to $l$ extended by the next term of the sequence.
\end{proof}

\begin{proposition}\label{p:barhilbert}
\begin{enumerate}
\item A discrete field is inductively Noetherian.
\item If $A$ is inductively Noetherian, so are $A/\mathfrak{a}$ and $S^{-1}A$.
\item $\Q[Y_1,\dots,Y_N]$ is inductively Noetherian~\cite{CP}.
\end{enumerate}
\end{proposition}

\begin{proof}[Sketch]
1. For a list $l$ in a discrete field, decide whether $\langle l\rangle$ is $0$ or $(1)$. In the
second case every extension is good. In the first case an extension $a$ is $0$ (good) or a unit,
after which every extension is good.
2. Induction on $\mathrm{Bar}$. At each use of rule~2, lift \emph{one} element.
3. Dickson's lemma is proved as an inductive bar on monomials, and Gröbner reduction transports it
to ideals.
\end{proof}

\begin{comment}
This is why the inductive form is used. A sequence in $A/\mathfrak{a}$ or $S^{-1}A$ cannot be lifted
to a sequence in $A$ without countable choice. In the inductive form only one element is lifted at a
time. Proposition~\ref{p:hilbert} follows from Proposition~\ref{p:barhilbert}, Lemma~\ref{l:seq}, and
the fact that a Rees algebra is a quotient of a polynomial ring.
\end{comment}

\begin{thebibliography}{9}
\bibitem{ATW} D.~Abramovich, M.~Temkin and J.~Włodarczyk, \emph{Functorial embedded resolution via
  weighted blowings up}, Algebra Number Theory 18 (2024), 1557--1587.
\bibitem{Brais} M.~J.-L.~Brais, \emph{Streamlining resolution of singularities with weighted
  blow-ups}, arXiv:2512.01859 (2025).
\bibitem{CP} T.~Coquand and H.~Persson, \emph{Gröbner bases in type theory}, in Types for Proofs and
  Programs, LNCS 1657 (1999), 33--46.
\bibitem{HM} C.~Hägg and A.~Mörtberg, \emph{A Bézout domain that is not an elementary divisor domain},
  arXiv:2609.35229, 2026, \url{https://arxiv.org/abs/2609.35229}. The proof was found with GPT-6 Astra.
  Lean formalization: \url{https://github.com/Zelaron/bezout-counterexample-lean}.
\bibitem{LQ} H.~Lombardi and C.~Quitté, \emph{Commutative Algebra: Constructive Methods},
  Springer, 2015.
\bibitem{Pm} H.~Perdry, \emph{An elementary proof of Krull's intersection theorem}, Amer.\ Math.\
  Monthly 111 (2004), 356--357.
\bibitem{P} H.~Perdry, \emph{Strongly Noetherian rings and constructive ideal theory},
  J.~Symbolic Comput.\ 37 (2004).
\end{thebibliography}

\end{document}
